\documentclass[11pt]{article}
\usepackage[margin=1in]{geometry}
\usepackage[T1]{fontenc}
\usepackage[utf8]{inputenc}
\usepackage{amsmath,amssymb,amsthm}
\usepackage{enumitem}

\title{ICPC World Finals 2025\\C. Bride of Pipe Stream}
\author{}
\date{}

\begin{document}
\maketitle

\section*{Problem Summary}

At every station we may split the incoming flow arbitrarily among the outgoing ducts, while each duct
then distributes its flow to lower stations and reservoirs by fixed percentages. The graph is acyclic
because every duct only goes to larger-numbered objects. We must maximize the minimum percentage of
the total source flow that reaches any reservoir.

\section*{Key Observations}

\begin{itemize}[leftmargin=*]
    \item Let $x=(x_1,\dots,x_r)$ be the vector of reservoir inflows produced by some station-splitting
    strategy. The set of all such vectors is convex, because every station can mix two strategies by using
    the corresponding convex combination of its split ratios.
    \item For any weight vector $\lambda \ge 0$ with $\sum \lambda_i = 1$, maximizing
    $\lambda \cdot x$ is easy. If $H_i(\lambda)$ denotes the best weighted value obtainable from station $i$,
    then
    \[
    H_i(\lambda)=\max_{\text{duct }d\text{ out of }i}
    \left(\sum_{j=1}^{r} \lambda_j \cdot \text{res}_{d,j}
    + \sum_{k} \text{frac}_{d,k} \cdot H_k(\lambda)\right).
    \]
    \item The original objective is
    \[
    \max_x \min_j x_j.
    \]
    Over a convex feasible set this is equal to
    \[
    \min_{\lambda \ge 0,\ \sum \lambda_i=1} \max_x \lambda \cdot x.
    \]
    So we only need to minimize $H_1(\lambda)$ over the simplex.
    \item Since $r \le 3$, the simplex has dimension at most $2$. The function $H_1$ is convex because it
    is the pointwise maximum of linear functions, so one-dimensional and nested golden-section searches
    are sufficient.
\end{itemize}

\section*{Algorithm}

\begin{enumerate}[leftmargin=*]
    \item Store every duct of every station as:
    the fixed fractions going to later stations, and the fixed fractions going directly to the reservoirs.
    \item For a fixed $\lambda$, evaluate $H_i(\lambda)$ by dynamic programming from station $s$ down to
    station $1$, since all edges go to larger indices.
    \item If $r=1$, the answer is just $H_1(1)$.
    \item If $r=2$, write $\lambda=(t,1-t)$ and minimize the convex function $H_1(t,1-t)$ on $[0,1]$ by
    golden-section search.
    \item If $r=3$, write $\lambda=(x,y,1-x-y)$ with $x,y \ge 0$ and $x+y \le 1$.
    For a fixed $y$, minimize over $x \in [0,1-y]$ by golden search, then minimize the resulting convex
    function of $y$ by another golden search.
\end{enumerate}

\section*{Correctness Proof}

We prove that the algorithm returns the correct answer.

\paragraph{Lemma 1.}
For a fixed weight vector $\lambda$, the dynamic program computes the maximum possible value of
$\lambda \cdot x$ obtainable from each station.

\paragraph{Proof.}
Consider a station $i$. Once the incoming flow reaches $i$, the only decision is how to split that flow
among its outgoing ducts. Because the weighted objective is linear, sending all flow through the outgoing
duct with the largest downstream weighted value is optimal. The formula used by the program is exactly
that best value. Processing stations in reverse order is valid because every duct goes only to larger
indices, so all needed subproblems are already known. \qed

\paragraph{Lemma 2.}
The optimum minimum reservoir inflow equals
\[
\min_{\lambda \ge 0,\ \sum \lambda_i=1} H_1(\lambda).
\]

\paragraph{Proof.}
Let $F$ be the convex set of all feasible reservoir vectors. For any feasible $x$ and any simplex weight
vector $\lambda$, we have $\min_j x_j \le \lambda \cdot x$. Therefore
\[
\max_{x \in F} \min_j x_j \le \min_{\lambda} \max_{x \in F} \lambda \cdot x.
\]
Conversely, if $z$ is the optimal max-min value, then the set $F$ intersects the box
$[z,\infty)^r$ and does not intersect the interior of any strictly larger such box. A supporting
hyperplane of $F$ at an optimal point gives a simplex weight vector $\lambda$ for which
$\max_{x \in F}\lambda \cdot x = z$. Hence equality holds. \qed

\paragraph{Theorem.}
The algorithm outputs the correct answer for every valid input.

\paragraph{Proof.}
By Lemma 1, for every tested $\lambda$ the program evaluates $H_1(\lambda)$ correctly. By Lemma 2,
the true answer is the minimum of this function over the simplex. Because $H_1$ is convex, its
restriction to any line segment is also convex, hence unimodal. Therefore the golden-section searches
used by the program converge to the global minimum in the one-dimensional and nested two-dimensional
cases covered by $r \le 3$. Thus the printed value is the optimal minimum reservoir inflow. \qed

\section*{Complexity Analysis}

Let $D$ be the number of ducts. One evaluation of $H_1(\lambda)$ costs $O(D \cdot r)$, which is
$O(D)$ because $r \le 3$. The search performs only a fixed number of evaluations, so the total running
time is $O(D)$ up to a constant factor, and the memory usage is $O(D)$.

\section*{Implementation Notes}

\begin{itemize}[leftmargin=*]
    \item The output is a percentage, so the program multiplies the final fraction by $100$.
    \item Using about $80$--$200$ golden-search iterations is easily enough for the required $10^{-6}$
    absolute error.
\end{itemize}

\end{document}
